\section{Latency under load}
\label{sec:latency}

A service cost implies a capacity, and a capacity implies a knee: below it delay
is service time, above it delay is queue and grows until something drops. The
offered rate is swept across that knee.

The load must be open-loop. A closed-loop generator waits for the previous
packet to be served before sending the next, so it cannot offer more than the
receiver's capacity, and the queueing this experiment exists to expose never
forms. Reporting a closed-loop curve as a latency-under-load result is a
standard error, and it always produces a flatteringly flat tail.

\noindent\textbf{What was measured, and on which clock.}
In the two-host topology, latency is a \emph{round trip} measured entirely on
the generator's own clock. The generator stamps the packet, the data plane
forwards it back, and the generator reads the stamp out of the header, which
travels in the clear as associated data and so needs no decryption. No clock
synchronisation between the hosts is assumed and no one-way delay is claimed.
The no-crypto echo baseline, measured the same way, is what makes the difference
interpretable. This run used the \latTopology\ topology under network condition
\latNetem, with the generator on a separate host: \latInterpretation.

\noindent\textbf{What this run produced.}
The capacity used to set the rate grid is measured on the same configuration
immediately beforehand rather than taken as a constant, so that the grid
straddles the knee on whatever machine this is. Here that capacity was
\latCapacityPPS\ packets per second. The rate sweep that follows it returned no
packets over the wire, so the knee is reported as \latKneePPS\ packets per
second, or \latKneeFraction\ of capacity, and no curve is drawn. That is a
recorded absence rather than a measured flat tail: the harness marks every point
of the sweep unavailable with its reason, and the analysis declines to read a
knee from it. Section~\ref{sec:threats} states what the missing curve costs the
rest of the argument.
